\section{Registered computational study}
\label{sec:computational-study}
\label{sec:experiments}

The registered algorithmic compression benchmark is a finite synthetic study of
tagged identity-closure instances. It can test exact-method agreement, certified
burden quality, preprocessing prevalence, and the declared adversarial
mechanisms. It cannot establish average-case complexity, steady-state runtime
rankings, or a population law for financial libraries.

\subsection{Registered design, execution, and audit}

The final registry contains \AORBenchmarkInstances\ analysis instances:
\AORBenchmarkSmallExact\ small exact, \AORBenchmarkStructured\ structured,
and \AORBenchmarkAdversarial\ adversarial. Family B requested variation in
strategy and tagged-row counts, coverage density, module overlap, bundle
prevalence, unique-carrier frequency, weight dispersion, and
frontier--module association. Family C contains the predeclared gap, tie,
duplicate, dominance, rare-carrier, symmetry, and bundle mechanisms. Failed
generation remains a terminal result; no failed row is silently repaired or
dropped.

The algorithm grid contains enumeration, requirement-mask DP, two registered
HiGHS formulations, weighted and cardinality greedy, greedy plus reverse
deletion, rechecked deletion orders, seeded random deletion, and 32-start
deletion. Maximum immediate additive burden release is exactly the
heaviest-safe-first ordering. The two labels denote one method in all scientific comparisons.

The audit passed all \AORBenchmarkRuns\ registered run units and accepted
\AORBenchmarkAccepted\ returned libraries after original-coordinate mandatory,
coverage, frontier, identity-closure, and exact-burden checks. It retained
\AORBenchmarkUnsuccessful\ unsuccessful records in their terminal categories.
HiGHS \texttt{OPTIMAL} remains solver evidence plus an exact incumbent check,
not formal or exhaustive proof.

\subsection{Exact agreement and applicability}

Among \AORExactCapable\ registry instances, all required exact-method burdens
were available and equal on \AORExactAgreement; \AORExactUnavailable\ did not
have a complete comparison. \AORImplementationErrors\ adversarial enumeration
jobs are recorded as
\texttt{IMPLEMENTATION\_ERROR}, but their substantive cause is known: the
residual optional count exceeded the declared enumeration cap. The recorded status is retained in the failure census. They are enumeration-inapplicable under that cap,
and exact agreement is unavailable rather than inferred. Optimizer identities
differed at equal burden on \AOROptimizerIdentityDifferences\ complete
comparisons, which records valid tied optima rather than disagreement.

These are exact finite facts about specified instances. The medium structured
references are HiGHS-optimal incumbents that pass exact feasibility and burden
rechecks; no second solver was pinned.

\subsection{Certified burden quality}

Table~\ref{tab:aor-heuristic-quality-compact} separates Family A exact finite
references from Family B exactly rechecked HiGHS-optimal references. Within
those distinct evidence classes, weighted greedy plus reverse deletion and
heaviest-safe-first often attain the available reference, while several other
deletion orders have larger observed tails. Multi-start attains every available
Family A reference but not every Family B reference. These are descriptive
outcomes of the registered instances, not approximation guarantees or an
algorithm ranking beyond the design.

\AORHeuristicQualityTable

The registered heaviest-safe-first gap family gives a different kind of
evidence. Its exactly rechecked finite ratios match
Theorem~\ref{thm:aor-heaviest-safe-first-gap} at every registered scale. The
theorem establishes the universal family statement; the experiment checks the
implementation on the named fixtures.

The supporting benchmark appendix reports the finite ratios.

\subsection{Preprocessing and realized design limits}

Full fixed-point preprocessing has a strong realized effect. In every
structured size cell, median variable and requirement reduction is complete
among available records. This is evidence that forced-carrier propagation is
important for these generators. It is simultaneously a limitation: most
structured instances do not exercise a difficult residual MIP or DP search, so
the study cannot identify general residual-search scaling.

\AORPreprocessingMultistartTable

The registered generator also contains infeasible requested cells. All
\AORGenerationFailed\ run units arising from
\AORGenerationFailedInstances\ structured source-instance cells are retained,
and the artifact audit shows that
the failures remove complete combinations of low density and bundle/unique-
carrier settings rather than occurring at random. Effects and interactions
requiring those absent cells are not estimable from this design.

Multi-start deletion completes all registered starts on
\AORMultistartComplete\ instances and improves on the first start on
\AORMultistartImproved. Memory-limited primary runs remain failures; successful
sensitivity records are not substituted into their denominators.

\subsection{Limits of the comparison}

Recorded times include first-call Julia compilation and contention among eight
one-thread workers on the recorded worktree. They do not support steady-state
speed rankings. The registered performance profile also mixes feasible-return
and completed-search targets, so it is not comparative evidence here.
Registered model outputs do not support structural inference: their HC3
uncertainty ignores repeated observations within instances, and entire
requested generator cells are unavailable. Neither model refitting nor missing
outcomes are used to improve the reported comparison.

Every unsuccessful run remains in the full denominator, including generation
failures, memory limits, interruptions, absent primals, and nonbinary returned
vectors. No run ended at the time limit. Online Resource~1, Section~S7 gives
the design, complete aggregate failure census, and exact-agreement audit.
